\chapter*{Abstract}

\begin{otherlanguage}{english}
\noindent Small and medium-sized enterprises (SMEs) increasingly need web
applications of their own, but they can rarely afford to build from scratch
the components that almost every SaaS (\textit{Software as a Service})
application repeats: secure authentication, management of users and teams
with roles, subscription billing, and a tool to administer data without
programming. This Bachelor's Thesis designs, implements and validates
PymeKit, a reusable software architecture that brings those components
together into a common base on which new SaaS applications for SMEs can be
developed in an agile way.

PymeKit combines a multi-tenant SaaS application and a content management
system (CMS) for the database into a single product. The contributions of
this work are the integration of the CMS into the administration console of
the application, with a single service, a single sign-in and a single library
stack; the reorganisation of the console into business areas that an SME can
understand; a user interface entirely in Spanish on top of an extensible
internationalisation infrastructure; the centralisation of the reusable
configuration; and the documented correction of numerous security flaws
found in adversarial audits.

The solution is organised as a monorepo of decoupled packages written in
strict TypeScript. The web application uses TanStack Start and React; data,
authentication and storage rely on Supabase (PostgreSQL); payments, on
Stripe; and the CMS API, built with Hono and Drizzle, is mounted inside the
web server itself. Authorisation is enforced in the database through
row-level security (RLS), so tenant isolation does not depend on the
application code alone. Access to the console requires a second
authentication factor, and every action performed from the CMS is recorded in
an audit log that staff cannot tamper with.

Quality is verified with unit tests, more than two thousand database tests
with pgTAP ---including cross-tenant attack tests--- and end-to-end tests with
Playwright, all run in continuous integration. More than fifty issues were
documented in a log during the process, many of them security flaws detected
in the audits and fixed during the integration. This report describes the
planning, analysis, design, construction, testing and deployment of the
platform, the evaluation of its reusability, and the developer, installation
and user manuals.

\vspace{.5cm}

\noindent \textbf{Keywords:} SaaS, SMEs, reusable architecture,
multi-tenancy, row-level security (RLS), Supabase, PostgreSQL, CMS, TanStack
Start.
\end{otherlanguage}
